\section{宣统转型可跳转事件节点}
以下锚点将咨议、铁路、军队、省级政治、谈判与退位分列，避免以单一革命中心或诏书取代各地不同的政治时间。
\subsection{1900—1909}
\eventanchor{EQ-1909-PROVINCIAL-ASSEMBLIES}{各省咨议局开办，地方精英获得新的…}宣统元年（1909年），各省咨议局开办，地方精英获得新的公开政治场域。核心取证：《宣统政纪》宣统元年相关条；宫中朱批奏折、内阁大库题本与《清史稿》相关志传。此类文件确认机构作出或收到的行政动作，不能跳过地方执行与反馈。来源保留：以宣统元年（1909年）的同期文书为定位起点；官修记录、奏报或外交文本服务特定机构立场，具体日期、规模、地方执行与对手视角须以独立材料复核。
\eventanchor{EQ-1909-RAILWAY-RECOVERY-MOVEMENT}{铁路国有化和路权问题引发各省商绅…}宣统元年（1909年），铁路国有化和路权问题引发各省商绅、公司和官府的争论。核心取证：《宣统政纪》宣统元年相关条；户部奏销、盐政、河工或海关档案。此类文件确认机构作出或收到的行政动作，不能跳过地方执行与反馈。来源保留：以宣统元年（1909年）的同期文书为定位起点；官修记录、奏报或外交文本服务特定机构立场，具体日期、规模、地方执行与对手视角须以独立材料复核。
\subsection{1910—1919}
\eventanchor{EQ-1910-ZIZHENG-YUAN}{资政院开院，中央层面的立宪咨询机…}宣统二年（1910年），资政院开院，中央层面的立宪咨询机构开始运作。核心取证：《宣统政纪》宣统二年相关条；宫中朱批奏折、内阁大库题本与《清史稿》相关志传。此类文件确认机构作出或收到的行政动作，不能跳过地方执行与反馈。来源保留：以宣统二年（1910年）的同期文书为定位起点；官修记录、奏报或外交文本服务特定机构立场，具体日期、规模、地方执行与对手视角须以独立材料复核。
\eventanchor{EQ-1910-FAMINE-AND-UNREST}{灾荒、物价和财政压力持续影响晚清…}宣统二年（1910年），灾荒、物价和财政压力持续影响晚清地方社会与政治动员。核心取证：《宣统政纪》宣统二年相关条；地方志、督抚奏折及地方档案。编年证词定位国家叙述中的时间，却需同地方或对方材料校验范围。来源保留：以宣统二年（1910年）的同期文书为定位起点；官修记录、奏报或外交文本服务特定机构立场，具体日期、规模、地方执行与对手视角须以独立材料复核。
\eventanchor{EQ-1911-RAILWAY-NATIONALIZATION}{清廷宣布铁路干线国有化并以外债筹…}宣统三年（1911年），清廷宣布铁路干线国有化并以外债筹资，保路运动迅速扩大。核心取证：《宣统政纪》宣统三年相关条；宫中朱批奏折、内阁大库题本与《清史稿》相关志传。此类文件确认机构作出或收到的行政动作，不能跳过地方执行与反馈。来源保留：以宣统三年（1911年）的同期文书为定位起点；官修记录、奏报或外交文本服务特定机构立场，具体日期、规模、地方执行与对手视角须以独立材料复核。
\eventanchor{EQ-1911-SICHUAN-RAILWAY-PROTEST}{四川保路运动升级，赵尔丰处置引发…}宣统三年（1911年），四川保路运动升级，赵尔丰处置引发更大政治危机。核心取证：《宣统政纪》宣统三年相关条；地方志、督抚奏折及地方档案。编年证词定位国家叙述中的时间，却需同地方或对方材料校验范围。来源保留：以宣统三年（1911年）的同期文书为定位起点；官修记录、奏报或外交文本服务特定机构立场，具体日期、规模、地方执行与对手视角须以独立材料复核。
\eventanchor{EQ-1911-WUCHANG-UPRISING}{武昌起义爆发，新军与革命团体控制…}宣统三年（1911年），武昌起义爆发，新军与革命团体控制武汉部分地区。核心取证：《宣统政纪》宣统三年相关条；军机处档、战事方略及地方奏报。编年证词定位国家叙述中的时间，却需同地方或对方材料校验范围。来源保留：以宣统三年（1911年）的同期文书为定位起点；官修记录、奏报或外交文本服务特定机构立场，具体日期、规模、地方执行与对手视角须以独立材料复核。
\eventanchor{EQ-1911-PROVINCIAL-INDEPENDENCE}{多省宣布脱离清廷，帝国行政和财政…}宣统三年（1911年），多省宣布脱离清廷，帝国行政和财政网络快速分裂。核心取证：《宣统政纪》宣统三年相关条；宫中朱批奏折、内阁大库题本与《清史稿》相关志传。此类文件确认机构作出或收到的行政动作，不能跳过地方执行与反馈。来源保留：以宣统三年（1911年）的同期文书为定位起点；官修记录、奏报或外交文本服务特定机构立场，具体日期、规模、地方执行与对手视角须以独立材料复核。
\eventanchor{EQ-1911-YUAN-SHIKAI-RETURN}{清廷起用袁世凯处理北方军政和议和…}宣统三年（1911年），清廷起用袁世凯处理北方军政和议和事务。核心取证：《宣统政纪》宣统三年相关条；宫中朱批奏折、内阁大库题本与《清史稿》相关志传。此类文件确认机构作出或收到的行政动作，不能跳过地方执行与反馈。来源保留：以宣统三年（1911年）的同期文书为定位起点；官修记录、奏报或外交文本服务特定机构立场，具体日期、规模、地方执行与对手视角须以独立材料复核。
\eventanchor{EQ-1912-REPUBLIC-FOUNDING}{中华民国临时政府在南京成立，与清…}宣统四年（1912年），中华民国临时政府在南京成立，与清廷退位谈判并行。核心取证：《宣统政纪》宣统四年相关条；宫中朱批奏折、内阁大库题本与《清史稿》相关志传。此类文件确认机构作出或收到的行政动作，不能跳过地方执行与反馈。来源保留：以宣统四年（1912年）的同期文书为定位起点；官修记录、奏报或外交文本服务特定机构立场，具体日期、规模、地方执行与对手视角须以独立材料复核。
\eventanchor{EQ-1912-QING-ABDICATION}{宣统帝颁布退位诏书，清帝国终结}宣统四年（1912年），宣统帝颁布退位诏书，清帝国终结。核心取证：《宣统政纪》宣统四年相关条；宫中朱批奏折、内阁大库题本与《清史稿》相关志传。此类文件确认机构作出或收到的行政动作，不能跳过地方执行与反馈。来源保留：以宣统四年（1912年）的同期文书为定位起点；官修记录、奏报或外交文本服务特定机构立场，具体日期、规模、地方执行与对手视角须以独立材料复核。
